% GENERATED, edit the French unit or the translation table
\mxnewpage\section*{EV4-B -- EAM preparation}\label{EV4-B}
\addcontentsline{toc}{section}{EV4-B -- EAM preparation}
\textbf{Paper reserved for the teacher}. Whole programme.\par Duration: 2 h. No calculator. Marked out of 20. Practice sessions in May and June; EAM date for the zone to be confirmed.
\subsection*{Part 1 -- Fluency (6 points)}
Twelve questions, \num{0,5} point each. Only one correct answer; no justification required. An incorrect or missing answer scores 0, with no loss of points.
\begin{enumerate}
\item Calculate $3+2\times3$.\par A. $9$\quad B. $12$\quad C. $15$\quad D. $8$
\item Calculate $\frac12+\frac13$.\par A. $\frac{5}{3}$\quad B. $\frac{5}{6}$\quad C. $\frac{2}{5}$\quad D. $\frac{1}{6}$
\item Calculate $\dfrac{2^3}{2^{1}}$.\par A. $8$\quad B. $2$\quad C. $4$\quad D. $1$
\item A value of 50 increases by \pc{20}. Final value?\par A. $70$\quad B. $40$\quad C. $60$\quad D. $10$
\item A value of 80 is obtained after a \pc{20} decrease. Initial value?\par A. $60$\quad B. $64$\quad C. $96$\quad D. $100$
\item A \pc{25} increase then a \pc{20} decrease: overall coefficient?\par A. $\frac{21}{20}$\quad B. $1$\quad C. $\frac{4}{5}$\quad D. $\frac{9}{10}$
\item The solution of $2x+3=0$ is\par A. $\frac{3}{2}$\quad B. $6$\quad C. $\frac{-3}{2}$\quad D. $-6$
\item The slope between $\pt{0;3}$ and $\pt{2;9}$ is\par A. $6$\quad B. $-3$\quad C. $3$\quad D. $2$
\item The mean of $3, 5, 7$ is\par A. $3$\quad B. $15$\quad C. $6$\quad D. $5$
\item A sector of $90^\circ$ represents what fraction of the total?\par A. $\frac{1}{3}$\quad B. $\frac{1}{2}$\quad C. $\frac{1}{4}$\quad D. $\frac{3}{4}$
\item A probability is $\dfrac{2}{5}$. Its complement is\par A. $\frac{3}{5}$\quad B. $\frac{1}{5}$\quad C. $\frac{2}{5}$\quad D. $\frac{5}{2}$
\item For $f(x)=x^2+3$, $f(-2)$ is\par A. $12$\quad B. $5$\quad C. $-1$\quad D. $7$
\end{enumerate}
\mxnewpage\subsection*{Part 2 -- Independent exercises (14 points)}
Every step must be justified. Exact results are expected unless stated otherwise.
\mxneed{12}
\subsection*{Exercise 1 -- An average price (5 points)}
A service offers several rates. 60 percent of customers choose plan A; among them, half pay 0 units and the other half pay 4 units. All the other customers pay 6 units. A customer is chosen at random and $X$ is the price they pay. The last question studies a uniform reduction, including for the rate that was initially zero.
\begin{enumerate}
\item Establish the probability distribution of $X$. \hfill(1 pt)
\item Calculate its expectation and variance. \hfill(2 pts)
\item All prices fall by 1 unit. Calculate the new mean and the new standard deviation. Interpret the negative price obtained for some customers. \hfill(2 pts)
\end{enumerate}
\mxneed{12}
\subsection*{Exercise 2 -- Exponential maximum (5 points)}
The function $f$ is defined on $\intff{0}{4}$ by $f(t)=t^2e^{-t}$. We look for its maximum on this interval. For the last comparison without a calculator, you may use the inequality given, $e^2>4$.
\begin{enumerate}
\item Calculate $f\prime$ in factorised form. \hfill(1 pt)
\item Determine the variation and the maximum. \hfill(2 pts)
\item Compare $f(4)$ exactly with the maximum, given $e^2>4$. \hfill(2 pts)
\end{enumerate}
\mxneed{12}
\subsection*{Exercise 3 -- Determining a position (4 points)}
In an orthonormal frame, the points $A\pt{0;0}$ and $B\pt{4;0}$ are fixed. The point $M$ moves on the line with equation $y=2$; its coordinates are therefore $M\pt{x;2}$, where $x$ ranges over $\R$.
\begin{enumerate}
\item Calculate $\vect{MA}\cdot\vect{MB}$. \hfill(1 pt)
\item Determine the values of x for which the two vectors are orthogonal. \hfill(1 pt)
\item Interpret the result using the circle with diameter AB. \hfill(2 pts)
\end{enumerate}
